% DFS 拓扑序：菱形 1->2, 1->3, 2->4, 3->4，从1出发、出邻居递增。
% 发现/完成：1:1/8, 2:2/5, 4:3/4, 3:6/7；完成序 4,2,3,1，逆序 1,3,2,4。Kahn（FIFO）得 1,2,3,4。
\begin{tikzpicture}[font=\figurefont\fontsize{8}{10}\selectfont,>=Stealth,
  vtx/.style={circle,draw=FigureStroke,line width=.8pt,minimum size=6.2mm,inner sep=0pt,text=FigureInk,fill=white},
  tree/.style={draw=FigureBlue,line width=1.2pt,->},
  other/.style={draw=FigureMuted,line width=.8pt,dashed,->},
  arc/.style={draw=FigureStroke,line width=.8pt,->},
  tm/.style={font=\figurefont\fontsize{7.5}{9}\selectfont,text=FigureInk},
  note/.style={font=\FigureNoteFont\fontsize{7.5}{9}\selectfont,text=FigureMuted}]
% (a)
\begin{scope}
  \node[vtx] (v1) at (0,0) {1};
  \node[vtx] (v2) at (1.2,0.8) {2};
  \node[vtx] (v3) at (1.2,-0.8) {3};
  \node[vtx] (v4) at (2.4,0) {4};
  \draw[tree] (v1)--(v2); \draw[tree] (v2)--(v4); \draw[tree] (v1)--(v3); \draw[other] (v3)--(v4);
  \node[tm,left=1pt of v1] {1/8};
  \node[tm,above=1pt of v2] {2/5};
  \node[tm,below=1pt of v3] {6/7};
  \node[tm,right=1pt of v4] {3/4};
  \node at (1.2,-2.55) {(a) 发现/完成时刻};
\end{scope}
% (b)
\begin{scope}[xshift=4.2cm]
  \def\dx{0.95}
  \tikzset{adj/.style={arc,bend left=50},skip/.style={arc,out=70,in=110,looseness=1.1}}
  % Kahn 次序 1,2,3,4
  \node[anchor=west] at (-0.3,1.5) {Kahn：1, 2, 3, 4};
  \foreach \v [count=\k from 0] in {1,2,3,4} \node[vtx,fill=FigureBlueFill] (k\v) at (\k*\dx,0.2) {\v};
  \draw[adj] (k1.north east) to (k2.north west); \draw[skip] (k1.north) to (k3.north);
  \draw[skip] (k2.north) to (k4.north); \draw[adj] (k3.north east) to (k4.north west);
  % DFS 逆完成序 1,3,2,4
  \node[anchor=west] at (-0.3,-0.48) {DFS 逆完成序：1, 3, 2, 4};
  \foreach \v [count=\k from 0] in {1,3,2,4} \node[vtx,fill=FigureGreenFill] (f\v) at (\k*\dx,-1.8) {\v};
  \draw[skip] (f1.north) to (f2.north); \draw[adj] (f1.north east) to (f3.north west);
  \draw[adj] (f2.north east) to (f4.north west); \draw[skip] (f3.north) to (f4.north);
  \node[align=center] at (1.42,-2.55) {(b) 两个合法拓扑序：每条弧都指向右方};
\end{scope}
\end{tikzpicture}
